% SHARD-ID: CA-27-GAUSSIAN-BOSON-RESIDUAL-CONDITIONS
% SHARD-TITLE: Lorentz Residual Conditions on Gaussian Coefficients
% SHARD-SUMMARY: Converts the boost-time and rotation residuals into explicit Taylor conditions on the scalar coefficient symbol.
% SHARD-SUMMARY: Shows why isotropic Klein-Gordon quadratic data are necessary but not sufficient for a one-particle relativistic scaling limit.
% SHARD-SUMMARY: Adds checked anisotropic and doubler counterexamples as compiler rejection witnesses.
% SHARD-KEYWORDS: Gaussian boson, Lorentz residual, Hessian, anisotropy, doubler, coefficient conditions

\section{Lorentz Residual Conditions on Gaussian Coefficients}
\label{sec:gaussian-boson-residual-conditions}

\subsection{Status, Sources, and Dependencies}

\textbf{Status: checked local coefficient calculus.}  The sourced free scalar
model supplies the target dispersion.  The general coefficient residuals are
local Taylor calculations from
\(\omega_\epsilon(k)^2=\sum_rV_re^{i\epsilon k\cdot r}\).  The examples in
this shard are checked in Julia.

Dependencies: CA-24--CA-26 and CONVENTIONS.md (j).

% Source: literature/md/2010.11121/2010.11121.md:614 -- physical mass convention for the lattice scalar.
% Source: literature/md/2010.11121/2010.11121.md:623 -- nearest-neighbour lattice Klein-Gordon dispersion.
% Source: literature/md/2010.11121/2010.11121.md:676 -- continuum Klein-Gordon dynamics.
% Source: references/cft/Schottenloher2008/Schottenloher2008.md:4186 -- Klein-Gordon operator in the free bosonic QFT example.
% Checked: test/runtests.jl, "Gaussian boson Lorentz Hessian examples".

\subsection{Gradient Residual}

For a scalar coefficient family,
\[
  \omega_\epsilon(k)^2=\sum_rV_re^{i\epsilon k\cdot r}.
\]
The boost-time residual from CA-26 is
\[
  R_a(k)
  =
  \frac{1}{2}\partial_a\omega_\epsilon(k)^2-c^2k_a.
\]
Explicitly,
\[
  \frac{1}{2}\partial_a\omega_\epsilon(k)^2
  =
  \frac{i\epsilon}{2}\sum_r r_aV_re^{i\epsilon k\cdot r}.
\]
When \(V_r=V_{-r}\), this is real and can be written as a sine sum:
\[
  \frac{1}{2}\partial_a\omega_\epsilon(k)^2
  =
  -\epsilon\sum_{r>0} r_aV_r\sin(\epsilon k\cdot r),
\]
where \(r>0\) denotes any fixed choice of one representative from each
\(\{r,-r\}\) pair.

\subsection{Low-Energy Hessian Condition}

At a symmetric minimum \(k=0\), the first derivative vanishes.  The Hessian is
\[
  \partial_a\partial_b\omega_\epsilon(0)^2
  =
  -\epsilon^2\sum_r r_ar_bV_r.
\]
Therefore the quadratic Lorentz-speed condition is
\[
  -\frac{\epsilon^2}{2}\sum_r r_ar_bV_r
  =
  c^2\delta_{ab}.
\]
This is the first finite-dimensional coefficient test.  In compiler form:
\[
  \mathsf{HessResidual}_{ab}
  :=
  -\frac{\epsilon^2}{2}\sum_r r_ar_bV_r-c^2\delta_{ab}.
\]
It must vanish at the chosen low-energy point before the model can be routed to
a Klein-Gordon continuum target.

\subsection{Nearest-Neighbour Klein-Gordon}

For
\[
  V_0=m^2+\frac{2d}{\epsilon^2},
  \qquad
  V_{\pm e_a}=-\frac{1}{\epsilon^2},
\]
one obtains
\[
  -\frac{\epsilon^2}{2}
  \sum_r r_ar_bV_r
  =
  \delta_{ab}.
\]
The gradient residual is the stronger finite-\(\epsilon\) expression
\[
  R_a(k)=\frac{\sin(\epsilon k_a)}{\epsilon}-k_a,
\]
which vanishes as \(O(\epsilon^2k_a^3)\) for fixed physical momentum.

\subsection{Anisotropic Counterexample}

Let the nearest-neighbour coefficients carry different speeds \(c_a\):
\[
  \omega_\epsilon(k)^2
  =
  m^2+\frac{2}{\epsilon^2}
  \sum_a c_a^2(1-\cos(\epsilon k_a)).
\]
Then
\[
  -\frac{\epsilon^2}{2}\sum_r r_ar_bV_r
  =
  c_a^2\delta_{ab}.
\]
Unless all \(c_a\) agree with the target speed, the boost-time residual cannot
match a single Lorentz cone.  The Julia test suite checks this as a deliberate
failure case, not as an approximate success.

\subsection{Doubler Counterexample}

The Hessian condition is necessary but not sufficient.  The symbol
\[
  \omega_\epsilon(k)^2
  =
  m^2+\frac{1}{\epsilon^2}\sum_a\sin^2(\epsilon k_a)
\]
has the correct quadratic Hessian at \(k=0\) when \(c=1\), but it has additional
low-energy minima at Brillouin-zone points with \(\sin(\epsilon k_a)=0\).
On even periodic grids this produces multiple minima.  Therefore the compiler
also needs a \emph{single-patch} or \emph{sector-selection} condition: a
continuum relativistic scalar target is attached only after choosing one
low-energy patch and proving the other minima decouple or are intentionally
retained as extra species.

\subsection{Compiler Rule}

For the scalar Gaussian tier, a Lorentz-target compile attempt must emit at
least:
\[
\begin{array}{ll}
\text{mass term} & \omega(0)^2,\\
\text{speed matrix} &
  -\frac{\epsilon^2}{2}\sum_r r_ar_bV_r,\\
\text{boost residual} & R_a(k),\\
\text{rotation residual} &
  k_b\partial_a\omega^2-k_a\partial_b\omega^2,\\
\text{boost-boost residual} & R_bX_a-R_aX_b,\\
\text{patch data} & \text{chosen minimum and exclusion of unwanted minima}.
\end{array}
\]
Only the mass and speed entries are finite Taylor data.  The residual and patch
entries are what prevent an anisotropic or doubler model from being mislabeled
as a single Klein-Gordon field.
